\medskip\noindent
\textbf{Solution types and junction set.}\\
\begin{itemize}
	\item \emph{Penalized region:}\quad
	$\displaystyle
	\mathcal{P}
	:= \bigl\{\,x:\text{$x$ has prefix }1^{(9/10)n}\bigr\}.
	$
	\item \emph{Deep‐Prefix Optima (DPO):} for $\rho\in(0.8,0.9)$,
	$\displaystyle
	\mathcal{D}_\rho
	:= \bigl\{\,x:\text{$x$ has prefix }0^L,\ |x|_1\ge \rho n\bigr\}.
	$
	\item \emph{Block‐Jump Region (BJR$_k$):} for integer $k\ge1$,
	$\displaystyle
	\mathcal{B}_k
	:= \bigl\{\,x:\text{$x$ has prefix }0^{2k+1}1^{(9/10)n-2k-1}\bigr\}.
	$
	\item \emph{Block‐Jump Optimum (BJO):}\quad
	the unique optimum in $\mathcal{B}_k$ is $y^\star = 0^{2k+1}1^{\,n-2k-1}$.
	\item \emph{Junction set:}\quad
	$\displaystyle
	J \;=\; \bigl\{\,x:\mathrm{LO}(x)=(9/10)n-1\bigr\}.
	$
\end{itemize}

\medskip\noindent
\textbf{Objective function (\textsc{LTJ}$_{k,\rho}$).}\\
With the above sets,
\[
\textsc{LTJ}_{k,\rho}(x)
=\begin{cases}
	0,                          & x\in\mathcal{P},\\[2pt]
	\mathrm{LO}\bigl(x_{i>2k+1}\bigr)+2k+1, & x\in\mathcal{B}_k,\\[2pt]
	n-1,                        & x\in\mathcal{D}_\rho,\\[2pt]
	\mathrm{LO}(x),             & \text{otherwise.}
\end{cases}
\]


Essentially the function consists of a  main \leadingones~ slope of increasing fitness until $(9/10)n -1$ leading ones are reached. From there it is necessary to add a constant number $2k$ of leading zeroes before the last $n/10$ leading ones can be inserted to reach the optimum. However, there is a trap consisting of many leading zeroes ($\log n$ ?) that has fitness value higher than any point other than the optimum. High mutation rates are likely to end up in the trap, while low mutation rates are likely to find the optimum. Naturally  if the values of the local and global optima are switched, then high mutation rates will be successful with high probability, while low mutation rates will fail. In the following we will show that the (1+1) Opt-IA with symmetric IPH efficiently optimises the function independent of the location of the optimum, while SBM will get trapped on the all-ones bit-string with high probability if the trap is globally optimal while static hypermutations with mutation potential get stuck in the trap with high probability when the all-ones bit string is optimal.


\medskip\noindent
\textbf{Target event.}\\
Let $\mathcal{E}_{\mathrm{junction}}$ denote the event that the algorithm reaches $J$ before
sampling any $x\in\mathcal{D}_\rho$ or any $x\in\mathcal{B}_k$.
%\end{definition}



\begin{theorem}[Expected fitness–evaluation cost in terms of \(\tau\)]
	\label{thm:main-tau-final}
	Let the block size \(k\ge 1\) be constant, \(r:=2k+1\), and set the IPH
	parameter to \(c=\tfrac13\).
	Let \(\tau=\tau(n)\ge n\) be the ageing threshold (in iterations).  
	For the $(1{+}1)$ Opt–IA with inversely proportional hypermutation and
	FCM($\ge$) acceptance, the expected number of fitness evaluations to optimise
	\textsc{LTJ}\(_{k,\rho}\) is
	
	\[
	\;
	\mathbb{E}[\text{\emph{evals}}]
	\;=\;
	O\!\bigl(
	n\,\tau
	\;+\;
	n^{\,2k+3}
	\;+\;
	n^{\,\beta+1}
	\bigr),
	\]
	
	\noindent
	where $\beta:=\log_{2}60\ <6$.
	
	
\end{theorem}
We first need to bound from below the probability that the algorithm visits the junction point after restarts in order to later argue that from that junction point both sets of optima will be accessed.

\begin{lemma}[Junction precedes DPO and $\mathcal{B}_k$ with probability $1-o(1)$]
	\label{lem:junction-first-density-Bk}
	Fix $c\in(0,\tfrac13)$ and a constant $k\ge 1$.
	Choose any $\rho\in(0.8,0.9)$ and any $\xi\in(0,\rho-\tfrac12]$ such that
	\[
	q^\star \;:=\; 2c\,\bigl(1-\rho+\xi\bigr) \;<\; e^{-2}.
	\]
	Starting from a uniformly random initial parent and with any fixed reference $x^\top$ 
	the event $\mathcal{E}_{\mathrm{junction}}$ satisfies
	\[
	\Pr\bigl(\mathcal{E}_{\mathrm{junction}}\bigr)\;=\;1-o(1).
	\]
\end{lemma}

\begin{proof}
	Let $T:=n^2$.
	
	\medskip\noindent
	\textbf{(A) Reaching $J$ quickly.}
	Before $J$ is hit, the first flip equals the leftmost zero with probability $1/n$ and is accepted
	immediately by FCM($\ge$).  If $X\sim\mathrm{Bin}(T,1/n)$ counts such increases then
	$\mathbb{E}[X]=n$ and a Chernoff bound gives
	\[
	\Pr\bigl(X<(9/10)n\bigr)\ \le\ e^{-a n}
	\]
	for some constant $a>0$. Hence $J$ is reached within $T$ iterations with probability $1-e^{-a n}$.
	
	\medskip\noindent
	\textbf{(B) DPO before $J$ is $o(1)$.}
	Write $\alpha:=|x|_1/n$ and $\alpha^\top:=|x^\top|_1/n$.
	
	\emph{Case B.1 (near threshold): $\alpha\ge \rho-\xi$ and $\alpha^\top\ge \rho-\xi$.}
	Using $|x\wedge x^\top|\ge |x|+|x^\top|-n$,
	\[
	\mathrm{HD}(x,x^\top)\ \le\ (2-\alpha-\alpha^\top)\,n\ \le\ 2(1-\rho+\xi)\,n,
	\]
	\[
	M(x)\ \le\ q^\star n\ <e^{-2}n
	.
	\]
	A DPO requires all $L$ leading bits among the $M$ flips; hence
	\[
	\Pr(\text{DPO in this iteration})
	\ \le\
	\Bigl(\tfrac{M}{n}\Bigr)^{L}
	\ \le\
	(q^\star)^{L}
	\ =\
	n^{-\log_{2}(1/q^\star)}.
	\]
	Because $q^\star<e^{-2}$, we have $\log_{2}(1/q^\star)>2$, so $T\cdot n^{-\log_{2}(1/q^\star)}=o(1)$.
	
	\emph{Case B.2 (far from threshold): $\alpha\le \rho-\xi$.}
	A DPO has $|x|_1\ge \rho n$, so at least $\xi n$ net $0\!\to\!1$ flips are required within at most
	$R\le M\le cn$ tail flips. With initial zero/one fractions giving negative drift,
	a Serfling (hypergeometric) tail bound yields
	\[
	\Pr\bigl(Z-O\ge \xi n\bigr)\ \le\ e^{-\gamma n}
	\]
	for some $\gamma>0$. Including the necessity to flip all $L$ prefix bits only decreases the probability.
	Summing over at most $T$ iterations remains $e^{-\gamma n}$.
	
	Therefore $\Pr(\text{DPO before }J)=o(1)$.
	
	\medskip\noindent
	\textbf{(C) Entering $\mathcal{B}_k$ before $J$ is $o(1)$ (using $k$ constant).}
	Let $r:=2k{+}1$ (a fixed constant).  For any parent $x\notin\mathcal{B}_k$:
	
	\emph{(C.1) Once $\mathrm{LO}(x)\ge r$.}
	Then the first $r$ bits of $x$ are $1$, so to enter $\mathcal{B}_k$ within an iteration, at the (FCM) acceptance time the offspring must differ from $x$ in at least the $r$ positions $\{1,\dots,r\}$ plus the next leftmost zero inside $S:=\{r{+}1,\dots,\tfrac{9}{10}n\}$, i.e., in at least $r{+}1$ specific coordinates. Thus, for some constant $C=C(k)$,
	\[
	\Pr\bigl(\text{enter } \mathcal{B}_k \text{ in one iteration}\bigr)
	\ \le\
	C\,\binom{n}{r+1}^{-1}
	\ =\
	\Theta\!\big(n^{-(r+1)}\big).
	\]
	Union-bounding over at most $T=n^2$ iterations gives a total contribution
	$O\!\big(n^{-(r-1)}\big)=O\!\big(n^{-2k}\big)=o(1)$.
	
	\emph{(C.2) While $\mathrm{LO}(x)< r$.}
	Let $T_0:=\lceil 4 n\log n\rceil$.  By Chernoff on the $\mathrm{Bin}(T_0,1/n)$ count of leftmost-zero
	flips, $\Pr(\text{reach }\mathrm{LO}\ge r\text{ within }T_0)\ge 1-e^{-a' n}$ for some $a'>0$.
	Condition on this event. Starting from a uniformly random parent, the number $v$ of zeros in
	$S:=\{r{+}1,\dots,\tfrac{9}{10}n\}$ satisfies $v=\Theta(n)$ w.h.p.; over the next $T_0$ iterations,
	accepted tail flips change $v$ by $\pm1$, and a bounded-differences (Azuma/Serfling) inequality keeps
	$v=\Omega(n)$ w.h.p.  Hence, throughout this window, any acceptance into $\mathcal{B}_k$ would require
	changing at least $s=\Omega(n)$ specified coordinates, giving per-iteration hazard
	$\binom{n}{s}^{-1}=e^{-a'' n}$ for some $a''>0$.  Union over $T_0=O(n\log n)$ iterations remains
	$e^{-a''' n}$.
	
	Combining (C.1) and (C.2), the probability to enter $\mathcal{B}_k$ at any time before $J$ is $o(1)$.
	
	\medskip\noindent
	\textbf{(D) Combine.}
	With probability $1-e^{-a n}$ we reach $J$ within $T$ steps by (A).  Subtracting the $o(1)$
	contributions from (B) and (C) yields
	\[
	\Pr\bigl(\mathcal{E}_{\mathrm{junction}}\bigr)
	\ \ge\
	\bigl(1-e^{-a n}\bigr)\ -\ o(1)\ -\ o(1)
	\ =\
	1-o(1).
	\]
\end{proof}



\begin{lemma}[Per-epoch probability of accepting the block-jump optimum]
	\label{lem:BJO-epoch}
	Fix a constant block length \(k\ge 1\) and write \(r:=2k+1\).
	Let the IPH constant be \(c=\tfrac13\) and let an ageing epoch have length
	\(\tau\ge n\).  
	If the epoch begins with the parent in the junction set
	\(J=\bigl\{x:\mathrm{LO}(x)=0.9n-1\bigr\}\),
	then the probability that the epoch accepts the unique block-jump optimum $y^{\star}=0^{\,2k+1}\,1^{\,n-2k-1}$ satisfies
	\begin{align*}
		p_B
		\;&:=\;
		\Pr\!\bigl(\text{$y^{\star}$ accepted during the epoch}\bigr)
		\;\\&=\;
		\Omega\!\Bigl(\min\!\bigl\{\,1,\;
		\tau\,n^{-(2k+2)}\bigr\}\Bigr),
	\end{align*}
\end{lemma}

\begin{proof}
	At the epoch start the reference \(x^\top\) is the best individual from the
	previous epoch.  
	The length-\(0.1n+1\) suffix of the new parent \(x\) is still uniformly random,
	so with probability at least \(\tfrac12\)
	\[
	\mathrm{HD}(x,x^\top)\;\ge\;\frac{n}{20}.
	\]
	Condition on this event, losing only a factor \(2\).
	Then the mutation potential in every subsequent
	iteration of the epoch satisfies
	\(
	M\ge c\,\tfrac{n}{20} = n/60.
	\)
	
	\paragraph*{Per-iteration success probability (revised, $M\ge r{+}1$).}
	Conditioned on $\mathrm{HD}(x,x^\top)\ge n/20$ we certainly have
	$M=\lfloor c\,\mathrm{HD}\rfloor\ge r+1$.
	Accepting $y^\star$ happens as soon as the \emph{first} $r{+}1$ distinct flips
	include \emph{exactly} the set
	\(\{1,\dots,r,\;0.9n\}\)
	in any order; if this occurs, FCM stops immediately and no other
	leading-one bit can be destroyed.
	
	Because flips are sampled uniformly without replacement, the probability of
	this event is
	\[
	p_{\text{iter}}
	=\prod_{i=0}^{r}\frac{r+1-i}{n-i}
	=\frac{(r+1)!}{(n)_{\,r+1}}
	=\Theta\!\bigl(n^{-(r+1)}\bigr)
	=\Theta\!\bigl(n^{-(2k+2)}\bigr),
	\]
	independent of the exact value of~$M$ (we only need $M\ge r{+}1$ to allow that
	many flips).
	
	\medskip\noindent
	\emph{Per-epoch success probability.}
	An epoch supplies \(\tau\) independent iterations, so a union bound yields
	\[
	p_B
	\;\ge\;
	1-(1-p_{\text{iter}})^{\tau}
	\;\ge\;
	\Omega\!\bigl(\min\{1,\;\tau\,n^{-(2k+2)}\}\bigr).
	\]
	
	\noindent
	The hidden constant depends only on the fixed block size \(k\).
\end{proof}

\begin{lemma}[Per-epoch probability of accepting a DPO]
	\label{lem:DPO-epoch}
	Fix the IPH constant \(c=\tfrac13\) and any density parameter \(\rho\in(0.8,0.9)\).  
	Let an ageing epoch have length \(\tau\ge n\) and begin with the parent in the
	junction set \(J\).  
	Then the probability that this epoch accepts \emph{some}
	deep-prefix optimum in \(\mathcal{D}_\rho\) satisfies
	\begin{align*}
		p_D
		\;&:=\;
		\Pr\!\bigl(\text{a string in }\mathcal{D}_\rho\text{ is accepted during the epoch}\bigr)
		\;\\&=\;
		\Omega\!\Bigl(\min\!\bigl\{\,1,\;
		\tau\,n^{-\,\beta}\bigr\}\Bigr),
	\end{align*}
	where \(\beta=\log_{2}60<6\). 
\end{lemma}
\begin{proof}[Proof of Lemma \ref{lem:DPO-epoch} (revised)]
	At epoch start the reference $x^\top$ is the best string from the previous
	epoch (from $J$ or $\mathcal{B}_k$), and the new parent $x$ itself lies in~$J$.
	Hence
	\[
	|x|_1\;=\;0.9n-1 \;+\; |x_{\text{suffix}}|_1.
	\]
	The suffix has length $0.1n+1$ and is uniform; its ones count is
	$\mathrm{Bin}(0.1n+1,\tfrac12)$ with expectation $0.05n+O(1)$.
	By a Chernoff bound
	\(
	|x_{\text{suffix}}|_1\ge 0.04n
	\)
	with probability $1-e^{-\Omega(n)}$.
	Thus, already before any mutation,
	\[
	|x|_1
	\;\ge\;
	0.9n-1 + 0.04n
	\;>\;
	0.94n
	\;>\;
	\rho n
	\qquad(\rho<0.9).
	\tag{A}
	\]
	
	\textbf{Step 1 – Linear mutation potential.}
	As in the block-jump argument we have
	\(
	\Pr\!\bigl(\mathrm{HD}(x,x^\top)\ge n/20\bigr)\ge\tfrac12
	\)
	(owing to the uniform suffix); condition on this event.
	Then \(M\ge n/60\) throughout the epoch.
	
	\textbf{Step 2 – Per-iteration flip of the \(L\) prefix bits.}
	Let \(L=\lfloor\log_2 n\rfloor\).
	A necessary part of generating a deep-prefix optimum is to flip the prescribed
	\(L\) prefix positions $\{1,\dots,L\}$ from $1$ to $0$ before FCM stops.
	Sampling without replacement,
	\begin{align*}
		p_{\text{prefix}}
		&=\Pr\bigl(\{1,\dots,L\}\subseteq\text{ first }L\text{ flips}\bigr)
		=\frac{(M)_{L}}{(n)_{L}}
		\;\ge\;
		(M/n)^{L}
		\;\\&\ge\;
		(1/60)^{L}
		=n^{-\beta},\qquad
		\beta:=\log_{2}60<6.
	\end{align*}
	
	\textbf{Step 3 – Density constraint is automatic.}
	When those $L$ flips occur, the total ones count in the string drops by
	\(\le L\le\log_2 n\).
	Using (A) and \(L=o(n)\),
	\[
	|x|_1-L\;\ge\;0.94n-\log_2 n
	\;\ge\;\rho n
	\quad\text{for all sufficiently large }n.
	\]
	Hence the density requirement \(|x|_1\ge\rho n\) is \emph{already satisfied}
	once the prefix flips have happened; no additional tail flips are needed, and
	FCM immediately accepts a DPO.
	
	\textbf{Step 4 – Per-epoch probability.}
	Each iteration satisfies \(p_{\text{iter}}\ge n^{-\beta}\).
	Across the \(\tau\ge n\) iterations of the epoch,
	\[
	p_D
	\;\ge\;
	1-(1-p_{\text{iter}})^{\tau}
	\;\ge\;
	\min\!\bigl\{1,\;\tau\,p_{\text{iter}}\bigr\}
	\;=\;
	\Omega\!\bigl(\min\{1,\;\tau\,n^{-\beta}\}\bigr),
	\]
	completing the proof.
\end{proof}

\textcolor{green}{Now that the lemmata are established we can prove our theorem.}
\begin{proof}
	\textbf{1.  Evaluation cost of one epoch.}
	Each iteration evaluates at most \(M(x)\le c\,n\;(=n/3)\) offspring, so an
	epoch of length \(\tau\) uses
	\[
	\text{eval}_{\text{epoch}}\le c\,n\,\tau=\Theta(n\,\tau).
	\]
	
	\textbf{2.  Per–epoch success probabilities.}
	Starting an epoch in the junction set \(J\) (prob.\ \(1-o(1)\)), the two
	lemmas give
	\[
	p_B=\Omega\!\bigl(\min\{1,\tau\,n^{-(2k+2)}\}\bigr),\quad
	p_D=\Omega\!\bigl(\min\{1,\tau\,n^{-\beta}\}\bigr),
	\]
	and \(B\) and \(D\) are mutually exclusive within one epoch.
	
	\textbf{3.  Expected number of epochs to see both optima.}
	A coupon–collector bound yields
	\[
	\mathbb{E}[T_{\mathrm{both}}]
	\;\le\;
	\frac1{p_B}+\frac1{p_D}
	\;=\;
	O\!\Bigl(
	\max\{1,n^{2k+2}/\tau\}
	+\max\{1,n^{\beta}/\tau\}
	\Bigr).
	\]
	
	\textbf{4.  Total expected evaluations.}
	Multiplying by the epoch cost,
	\begin{align*}
		\mathbb{E}[\text{evals}]
		&\le
		\Theta(n\,\tau)\;
		O\!\Bigl(
		\max\{1,n^{2k+2}/\tau\}
		+\max\{1,n^{\beta}/\tau\}
		\Bigr)
		\\&=
		O\!\bigl(
		n\,\tau
		+n^{2k+3}
		+n^{\beta+1}
		\bigr).
	\end{align*}
\end{proof}

Thus the runtime is polynomial in \(n\) for every fixed \(k\) and for any
\(\tau=\mathrm{poly}(n)\).  If \(\tau\le n^{2k+2}\) and
\(\tau\le n^{\beta}\) the leading term becomes
\(O\!\bigl(n^{2k+3}+n^{\beta+1}\bigr)\);  
if \(\tau\) is larger than both thresholds, the waiting cost \(n\,\tau\)
dominates.
